\chapter{Everything So Far Was Turing Complete}
\label{ch:ent}
% ===================================================================

The trick, as the last chapter told it, was a restriction of scope.
Confine attention to computation and an ontology arrives free, a grounded
hypothesis language arrives free, and a great many questions that looked
metaphysical turn out to have answers that are merely expensive.
The bill for that free lunch was located to a single line.

This chapter presents a second bill, and it is the one that governs the
rest of the act.
Restricting scope to computation is a restriction to \emph{Turing}
computation, and everything built in the Turn stayed inside it.

It is worth being blunt about this, because the machinery was elaborate
enough to disguise it.
The rho calculus is Turing complete.
Metering it does not add power --- the cost monad of
Chapter~\ref{ch:costmonad} takes a theory and returns a theory that does
strictly less, since a term which cannot pay does not fire.
Logging it does not add power either; the history monad records what
happened and computes nothing new.
Decorating rewrites with semiring values, imposing reversibility,
imposing arbitrage-freedom: all of these carve out subcategories, and a
subcategory of a category of Turing-complete theories contains only
Turing-complete theories.
Part~\ref{part:physics} descended a ladder and never once left the room.

And yet the learner of Part~\ref{part:ecology} does something that a
budget-constrained Turing machine does not obviously do.
It \emph{picks}.
When two of its sends contend for one receive, one of them wins, and
nothing in the term determines which.
This chapter is about that gap: where it is, what fills it, and why the
answer has the shape that Chapter~\ref{ch:homeo-ent} gave to entropy.

\section{The closed room}

Recall the two constructions of Part~\ref{part:machinery} and what each
does to expressive power.

The cost monad installs a meter at the interaction cut.
Its effect on the transition system is subtractive: transitions that were
available become unavailable when the stack is exhausted.
Chapter~\ref{ch:costmonad} makes this precise as graded adequacy --- the
metered theory is adequate for the unmetered one up to the grading, which
is to say it distinguishes no more and sometimes less.

The history monad reifies rewrite events as terms.
Its effect is additive but not powerful: a term can now name a step it
took, which is exactly what Chapter~\ref{sec:reversibility} needs to
build the time-symmetric envelope, but naming a step one has taken is
not computing anything one could not compute before.

So neither construction moves a theory out of its Turing degree, and the
axioms of Part~\ref{part:physics} are restrictions rather than
extensions.
The ladder of that part is a ladder within a single room.

\begin{remark}[Why this is not a defect to be patched]
\label{rmk:ent-not-a-defect}
The temptation at this point is to bolt on an oracle: declare that some
distinguished term has access to a halting decision, and proceed.
That would be a definition, not a discovery, and it would leave entirely
open the question of what such a term is doing physically.
The route taken here is the opposite one.
We look for a place in the machinery that is \emph{already}
underdetermined --- a place where the theory as stated does not say what
happens next --- and ask what would have to be supplied there.
The claim of these chapters is that there is such a place, that it is not
exotic, and that a cost-accounted concurrent calculus is full of them.
\end{remark}

\begin{remark}[Why this belongs in the Prestige and not the Turn]
\label{rmk:why-here}
An earlier arrangement of this book put everything that follows in the
Turn, as its fifth part.
It was moved here, and the move is not cosmetic.

The Turn's business is construction: build a knower, give it a world,
count what a population of knowers can find out.
All of that is done inside the restriction, and doing it inside the
restriction is exactly what makes it work --- the ontology and the
hypothesis language are free precisely because the scope was narrow.
What follows is not more construction.
It is the audit of the restriction: what was outside it, who would live
out there, and what it would cost to notice.
That is the Prestige's business, and putting it anywhere else disguised
the fact that the whole apparatus of the Turn has a ceiling.

The reader should also be told what this costs us.
The tower of Chapter~\ref{ch:tower} is the least-built construction in
the book, and Chapter~\ref{ch:tower}'s own closing section says so.
Moving it here puts the book's thinnest ice under its loudest claims.
We think that is the honest arrangement rather than the flattering one:
a reveal that hid its weakest step would be the kind of trick the last
chapter just finished disowning.
\end{remark}

\section{Choice is what fills the gap}

The place where a concurrent theory does not say what happens next is the
race.
When two sends contend for a single receive, the reduction relation
offers both continuations and the calculus is silent about which occurs.

The traditional reading is that this is harmless underspecification, to be
resolved by an implementation and forgotten.
The reading taken here, and developed in Chapters~\ref{ch:apr}
and~\ref{ch:agy}, is that the resolution is a mathematical object with a
measurable strength.
Picking a winner from each contending family is a choice function on that
family.
Asserting that a complete schedule exists --- that every race, for the
whole run, gets resolved --- is an instance of the axiom of choice, and
which instance depends on the size and structure of the families
involved.
Finite contention is free.
Contention indexed by $\omega$ requires countable or dependent choice.
Class-sized families require the full axiom.

So the extra power is not bolted on.
In any calculus where interaction is primitive and confluence fails, it
is already present at every cut, wearing the name \emph{scheduler}.
Chapter~\ref{ch:apr} distinguishes the cuts at which this is observable
--- the \emph{apertures} --- from the cuts at which it is not, which is
the substantive part of the analysis and the reason the chapter is called
what it is.
Chapter~\ref{ch:chs} then makes the correspondence a typing discipline:
choice principles are the types, and schedulers are the terms.

\begin{remark}[The learner had something to decide, and this is what it was]
\label{rmk:ent-handoff}
Chapter~\ref{ch:homeo-ent} closed by observing that a theory sitting
strictly between the two erasures --- knowing something about its own past
it cannot express as a cost, and something about its own costs it cannot
express as a history --- is a computation with something to decide.
That was left as a promissory note.
This is what it was a note for.
The information of that shape is not decision-\emph{relevant} in some
loose sense; it is what makes the race at a cut into an aperture rather
than a formality, and the strength of the choice principle required to
resolve it is what the rest of these chapters measures.
\end{remark}

\section{Entropy and choice as one dial}

We can now say what the title of Chapter~\ref{ch:homeo-ent} was pointing
at.

Entropy production, in the sense of Chapter~\ref{sec:cost}, measures how
much which-way information the dynamics discards per step.
Choice strength, in the sense of Chapter~\ref{ch:apr}, measures how much
must be supplied to say which way it went.
These are the same quantity approached from opposite sides, and the
conjecture that they are related quantitatively is the natural one.

\begin{conjecture}[Dissipation bounds choice]
\label{conj:dissipation-choice}
Let $c$ be a cut in a cost-accounted interactive theory, and let
$\theta(c)$ be the dissipation charged to the interaction at $c$.
Then the Weihrauch degree of the scheduling problem at $c$ is bounded
below by a function of $\theta(c)$: a cut that dissipates nothing carries
no observable choice, and a cut whose scheduling problem requires
dependent choice must dissipate at least the corresponding erasure cost.
\end{conjecture}

The forward half of this is close to available.
A frictionless cut is one whose enzyme is value-preserving in both
directions, and a cut which is reversible in that sense is not an
aperture by Definition~\ref{ndo:def:aperture}: there is nothing to
observe about which way it went, because it can be run backwards.
The converse half is open, and the obstruction is that the Weihrauch
lattice is not linearly ordered while the dissipation is a single real
number, so no function of $\theta$ alone can pin down a degree.
The honest statement is probably that $\theta$ bounds the degree from
below along each chain, which is weaker and may be all that is true.
Chapter~\ref{ch:chs} grades the correspondence in the Weihrauch lattice
and is the place to look for what is actually established.

\section{What this buys, and what it does not}

It does not buy hypercomputation for free.
Nothing here exhibits a term that decides the halting problem, and the
tower of Chapter~\ref{ch:tower} remains a construction whose inhabitants
are hypothesized rather than built.

What it buys is a location.
If a physical system is to exceed what a Turing machine can do, the
excess has to be spent somewhere, and this chapter says where: at the
apertures, in the resolution of races, paid for in dissipated free
energy.
That is a considerably more constrained claim than the usual ones, and it
is falsifiable in the direction that matters --- if it turned out that
every aperture in a physically realizable system carried only finite
contention, then the framework would predict that no physical system
exceeds the Turing computable, and the question would be settled the
other way.

There is also a consequence for Part~\ref{part:ecology} which
Part~\ref{part:origins} has already made quantitative, and which the
reader met there without being told what it was for.
A population of learners has apertures wherever two of its members
contend, and the number of such places grows with the structural
elaboration of the population rather than with its raw size.
A large uniform population has few apertures; a small deeply nested one
has many.
Whatever measures nesting therefore measures the population's capacity to
be more than the sum of its computations --- and
Corollary~\ref{cor:ai-choice} argues that the thing which measures
nesting is the assembly index.

\section{What follows}

Chapter~\ref{ch:tower} constructs the tower of oracle extensions over
Turing-complete GSLTs, indexed by ordinals --- which is how anyone first
gets the intuition and is not how the thing actually works.
Its closing section says, in the chapter itself, what it gets wrong.
Chapters~\ref{ch:apr} and~\ref{ch:agy} replace the ordinal chain with
what is underneath it: cuts, apertures, and the question of where an
agent's boundary falls.
Chapter~\ref{ch:chs} makes the correspondence a typing discipline, with
choice principles as the types and schedulers as the terms.
Chapter~\ref{sec:consciousness} then asks who lives at a position in that
lattice, and answers that consciousness is the position --- a coordinate
every agent has, rather than a club with an entrance exam.
Chapter~\ref{sec:hyper-gaps} catalogs what is owed.

Only then does the act turn to reference.
Chapter~\ref{ch:notation} asks what a symbol system would have to be;
Chapter~\ref{ch:sim} asks what the answer implies for the hypothesis that
this world is simulated; and Chapter~\ref{ch:kant} gathers the whole
argument into a metaphysics.
The ordering is deliberate.
One cannot ask what it would take to ground a symbol system without first
knowing whether the world doing the grounding is richer than the world
being grounded, and that is a question about position in a lattice.

\begin{remark}[A warning about two ladders]
\label{rmk:ladders-again}
The ladder in these chapters and the ladder of Part~\ref{part:physics}
are perpendicular, and confusing them is the easiest available mistake.
Remark~\ref{rmk:two-ladders} states the difference.
Moving down the physics ladder is imposing axioms: it carves out
subcategories of theories and never leaves the Turing degree.
Moving up here is acquiring choice strength: it changes what can be
resolved and therefore what exists.
A theory can be low on one and high on the other.
\end{remark}

% ===================================================================
